\documentclass[11pt,a4paper]{article}
\usepackage[margin=2.5cm]{geometry}\usepackage{pdfpages}\usepackage{amsmath}
\usepackage[bookmarks=true,bookmarksopen=true]{hyperref}\usepackage{bookmark}
\begin{document}
\pagestyle{empty}\pdfbookmark[0]{Cover}{cover}
\begin{center}\vspace*{3cm}
{\LARGE\bfseries Centre-of-Expansion Shell Cosmology:\\[4pt] Collected Draft Papers\par}\vspace{1.5cm}
{\Large Chris Alameddin\par}\vspace{0.5cm}{\large October 2026\par}\vspace{0.8cm}
{\large\bfseries DRAFT --- work in progress\par}\vspace{2.5cm}\end{center}
\noindent\small Full table of contents on the next page. Page numbers refer to this combined PDF; each paper keeps its own original page numbering.
\clearpage
\pdfbookmark[0]{Contents}{toc}
{\noindent\Large\bfseries Contents\par}\medskip
{\small
\noindent\textbf{1. CoE cosmology: Apparent acceleration from an off-centre finite geometry}\dotfill\textbf{3}\par\smallskip
\noindent\hspace*{1.2em}Introduction: the inverse problem of apparent acceleration\dotfill 3\par
\noindent\hspace*{1.2em}Off-centre geometric solution and prior angular predictions\dotfill 5\par
\noindent\hspace*{1.2em}Migkas comparison — out-of-sample angular morphology\dotfill 6\par
\noindent\hspace*{1.2em}Exact relativistic (LTB) calculation of $A/A_{\rm all}$\dotfill 14\par
\noindent\hspace*{1.2em}CMB null test — the purely geometric model\dotfill 25\par
\noindent\hspace*{1.2em}Physical completion — motion, gravity, and photon propagation\dotfill 27\par
\noindent\hspace*{1.2em}Residual multipoles and limits\dotfill 34\par
\noindent\hspace*{1.2em}Independent signed matter-dipole constraint\dotfill 37\par
\noindent\hspace*{1.2em}Redshift evolution of $A(z,\theta)$\dotfill 40\par
\noindent\hspace*{1.2em}Discussion and conclusions\dotfill 41\par
\noindent\hspace*{1.2em}References\dotfill 43\par
\noindent\hspace*{1.2em}Appendix figures\dotfill 44\par
\noindent\hspace*{1.2em}Acknowledgments\dotfill 48\par
\noindent\textbf{2. Reading the CMB fingerprint of the CoE shell model}\dotfill\textbf{49}\par\smallskip
\noindent\hspace*{1.2em}Introduction\dotfill 49\par
\noindent\hspace*{1.2em}The fingerprint\dotfill 49\par
\noindent\hspace*{1.2em}Method\dotfill 50\par
\noindent\hspace*{1.2em}Results\dotfill 52\par
\noindent\hspace*{1.2em}Comparison with observed structure\dotfill 66\par
\noindent\hspace*{1.2em}Editing filled webs to reach the CMB level\dotfill 68\par
\noindent\hspace*{1.2em}Mesh, area and energy of an edited web\dotfill 68\par
\noindent\hspace*{1.2em}Crystal analogy and the compensation rule\dotfill 71\par
\noindent\hspace*{1.2em}A compensated random web reaches the CMB level\dotfill 71\par
\noindent\hspace*{1.2em}Pair energies and forces\dotfill 75\par
\noindent\hspace*{1.2em}Concentric shell models\dotfill 75\par
\noindent\hspace*{1.2em}Caveats\dotfill 89\par
\noindent\hspace*{1.2em}Conclusions\dotfill 89\par
\noindent\hspace*{1.2em}References\dotfill 90\par
\noindent\textbf{3. Pattern formation (v5)}\dotfill\textbf{91}\par\smallskip
\noindent\hspace*{1.2em}Introduction\dotfill 91\par
\noindent\hspace*{1.2em}A common framework\dotfill 91\par
\noindent\hspace*{1.2em}Classification\dotfill 93\par
\noindent\hspace*{1.2em}Reference case: Chladni figures (forced mode, D)\dotfill 93\par
\noindent\hspace*{1.2em}Class A: bulk finite-wavelength instabilities\dotfill 94\par
\noindent\hspace*{1.2em}Class B: equilibrium modulated phases\dotfill 98\par
\noindent\hspace*{1.2em}Class C: coarsening\dotfill 99\par
\noindent\hspace*{1.2em}Class D: interfacial instabilities\dotfill 100\par
\noindent\hspace*{1.2em}Class E: fronts and waves\dotfill 102\par
\noindent\hspace*{1.2em}Self-similarity\dotfill 104\par
\noindent\hspace*{1.2em}Chaos\dotfill 104\par
\noindent\hspace*{1.2em}Summary: what the patterns share\dotfill 105\par
\noindent\hspace*{1.2em}Wrinkled thin shells: supernova remnants and the CoE onion\dotfill 106\par
\noindent\hspace*{1.2em}References\dotfill 121\par
\noindent\hspace*{1.2em}Summary Table\dotfill 123\par
\noindent\textbf{4. Does everything have a centre? A philosophical note}\dotfill\textbf{125}\par\smallskip
}
\clearpage
\includepdf[pages=-,addtotoc={1,section,1,{CoE cosmology: Apparent acceleration from an off-centre finite geometry},p1}]{/workspace/cosmo_preprint/coe_shell_preprint.pdf}
\includepdf[pages=-,addtotoc={1,section,1,{Reading the CMB fingerprint of the CoE shell model},p2}]{/workspace/clump_scan/report/brief/main.pdf}
\includepdf[pages=-,addtotoc={1,section,1,{Pattern formation (v5)},p3}]{/workspace/pattern_formation/v5/pattern_formation_v5.pdf}
\pagestyle{plain}
\phantomsection\pdfbookmark[0]{Does everything have a centre?}{phil}
\begin{center}{\Large\bfseries Does everything have a centre? A philosophical note\par}\medskip{\large Chris Alameddin\par}\end{center}
\section*{1. The thesis}
Every physical object or system that has ever been observed has a centre. Depending on what is being described, that centre may be geometric (the centroid of a shape), dynamical (the centre of mass), or kinematic (the centre-of-momentum frame, in which the total momentum vanishes). A galaxy, a cluster, a star, a supernova remnant, a planetary system and a laboratory sample all admit one or more of these. Even an irregular or diffuse distribution of matter has a well-defined first moment,
\begin{equation}
\mathbf{x}_{\rm c}=\frac{\int\rho(\mathbf{x})\,\mathbf{x}\,d^3x}{\int\rho(\mathbf{x})\,d^3x},
\end{equation}
provided the total mass is finite. The integral fails to define a centre only when the distribution is infinite in extent or exactly uniform without bound.

The second half of the thesis is empirical. Over many discussions, the author has asked for counter-examples, including from several AI systems. Every example proposed of something \emph{without} a centre has turned out to be a mathematical abstraction, and none of them has ever been observed by any scientist as a physical object.
\section*{2. The candidate counter-examples}
The usual candidates fall into a few classes:
\begin{itemize}
\item \textbf{Infinite uniform space, an infinite plane or an infinite line.} These are idealisations. No finite measurement can establish that something extends without bound, and any observed region of them is finite and so has a centroid.
\item \textbf{The surface of a sphere, taken as an intrinsic two-dimensional manifold.} No point on the surface is special. But every physical sphere (a balloon, a planet, a soap film) is embedded in three-dimensional space and has a centre there. A centreless sphere exists only if you deliberately forget the embedding.
\item \textbf{Fractal limits.} A physical fractal-like structure (a coastline, a turbulent cascade, galaxy clustering on intermediate scales) is cut off at small and large scales. The ideal infinite-scale fractal is never observed.
\item \textbf{The homogeneous, isotropic FLRW universe.} Under the cosmological principle every point is equivalent, so there is no centre. This is the case that matters here, and it is taken up next.
\end{itemize}
In each case the absence of a centre comes from an idealising step: infinite extent, exact uniformity, or ignoring the embedding. It does not come from a measurement.
\section*{3. ``No centre'' in cosmology is an assumption}
The standard model's statement that the universe has no centre is not an observational result. It follows from the cosmological principle, which assumes spatial homogeneity and isotropy and thereby selects the FLRW metrics~\cite{rindler,ellis79}. Two statements are often merged and should be kept apart:
\begin{enumerate}
\item ``The CMB is very nearly isotropic about us.'' This is data.
\item ``Therefore there is no centre anywhere.'' This is not a deduction from (1). Isotropy about one point plus the Copernican assumption that we are not special implies homogeneity, but that assumption is a philosophical preference, not a measurement~\cite{ellis79,ellis02}.
\end{enumerate}
Spherically symmetric inhomogeneous models such as Lema\^itre--Tolman--Bondi solutions, and the earlier kinematic cosmology of Milne~\cite{milne}, show that a universe with a centre is mathematically consistent with general relativity and with many of the same observations. Ellis and collaborators have stressed that what we see on large scales depends on how local structure is averaged and fitted to a smooth model~\cite{ellis02,ellisstoeger}. Whether a centre exists is therefore an empirical question that the cosmological principle settles by assumption.
\section*{4. Connection to the CoE model}
The papers in this collection take the opposite starting point. If the universe began as an expansion from a point, a finite expansion necessarily has a centre, the centre of expansion (CoE). We sit a distance $D$ from it. Nothing in this requires local speeds above $c$: bulk recession of a distant shell faster than $c$ is expansion, not local motion.

On this view the thesis of Section~1 applies to the universe itself. A universe without a centre would be the one physical system that behaves like the mathematical abstractions of Section~2, and would be the only such system ever claimed.
\section*{5. What could test it}
A centre is not just a philosophical position. An off-centre observer should see direction-dependent signatures:
\begin{itemize}
\item a dependence of the local Hubble rate on direction relative to the CoE, measured in the papers as $A/A_{\rm all}(\theta)$, the per-angle expansion rate divided by the full-sky mean, compared with cluster-scaling anisotropy results such as those of Migkas et al.~\cite{migkas};
\item dipoles that do not line up with the kinematic CMB dipole, such as the quasar number-count dipole reported by Secrest et al.~\cite{secrest};
\item a spread in the inferred $H_0$ with direction and distance, which bears on the Hubble tension.
\end{itemize}
These signals may also have conventional explanations: survey systematics, local bulk flows, or ordinary structure. A null result at sufficient precision would push the CoE far away or remove it. A consistent, aligned set of anisotropies would be evidence that the universe, like everything else observed, has a centre.
\section*{6. Balance}
None of this proves that the universe has a centre. The argument is an inductive one: every observed system has a centre, and every centreless example is an abstraction. Induction can fail, and the universe may be unique. But the burden is symmetric. ``No centre'' should be treated as a hypothesis to be tested, on the same footing as the CoE, and not as a settled result.
\begin{thebibliography}{9}
\bibitem{ellis79} G.~F.~R. Ellis, ``The homogeneity of the universe,'' \emph{Gen.\ Rel.\ Grav.} \textbf{11}, 281 (1979).
\bibitem{ellis02} G.~F.~R. Ellis, ``Cosmology and local physics,'' \emph{New Astron.\ Rev.} \textbf{46}, 645 (2002).
\bibitem{ellisstoeger} G.~F.~R. Ellis and W.~Stoeger, ``The `fitting problem' in cosmology,'' \emph{Class.\ Quantum Grav.} \textbf{4}, 1697 (1987).
\bibitem{milne} E.~A. Milne, \emph{Relativity, Gravitation and World-Structure} (Oxford, 1935).
\bibitem{rindler} W.~Rindler, \emph{Relativity: Special, General, and Cosmological}, 2nd ed.\ (Oxford, 2006).
\bibitem{migkas} K.~Migkas et al., ``Probing cosmic isotropy with a new X-ray galaxy cluster sample through the $L_X$--$T$ scaling relation,'' \emph{A\&A} \textbf{636}, A15 (2020).
\bibitem{secrest} N.~J. Secrest et al., ``A test of the cosmological principle with quasars,'' \emph{ApJL} \textbf{908}, L51 (2021).
\end{thebibliography}
\end{document}
